\chapter{Background and Related Work}

This chapter argues one line of reasoning end to end, across eight sections, rather than
surveying eight bodies of literature side by side. In order: is the problem real, what has
already been tried, has removing things instead of adding them been tried and did it work,
can any of this be built on hardware people actually own, what does a system built that way
need to be tested for, and how should that test be run.

Section 2.1 asks the first question. It establishes why peripheral events capture attention
involuntarily and why that capture has a cost, which is the intuition the rest of the chapter
builds on. Section 2.2 reviews the techniques already used to guide attention in mediated
views, from subtle gaze direction through saliency modulation to overt cues. Section 2.3 asks
whether removing things, rather than adding them, has been tried before: it reviews diminished
reality (DR) and the newer framing of visual noise cancellation, the conceptual home of this
work. Section 2.4 asks what a human periphery can tolerate, what it can lose without harm, and
where the comfort limits on restricting it further sit. Section 2.5 asks the hardware question
directly: what a consumer passthrough headset actually permits, what has already been built on
one, and reports a documented search establishing that this dissertation's own architecture has
no direct precedent. Section 2.6 surveys the applied contexts, driving, workstation, and
educational, that supply this dissertation's test scenarios, answering what to test for.
Section 2.7 reviews how XR attention systems are evaluated, answering how to test it, and
grounds the methodological choices of Chapter 6. Section 2.8 closes the chapter: it draws the
research gaps out of everything above and positions this work against its two nearest
published neighbours.

\section{Foundations: Attention, Salience, and Cognitive Load}

\subsection{Salience and attentional capture}

The computational account of visual attention that underpins this entire dissertation is the
salience-map model of Itti, Koch and Niebur (1998): attention is drawn to locations whose
centre–surround contrast, in intensity, colour opponency, or orientation, differs most from their
neighbourhood. Three consequences of this model recur throughout the design of the present system.
First, \textit{contrast is the master variable}: any manipulation that raises local contrast attracts gaze,
and any manipulation that flattens it repels gaze. Second, salience is channel-specific, and the
channels are not equal. In every head-to-head comparison in the guidance literature, luminance
manipulations outperform chromatic ones (Bailey et al., 2009; Grogorick et al., 2017; Waldin et al.,
2017). Third, salience operates \textit{pre-attentively}: capture by a high-contrast peripheral event
happens before, and regardless of, the observer's intentions. Barreiros, Veas and Pammer-Schindler
(2016) argue that AR visualisations must operate at exactly this pre-attentive level to feel automatic
rather than effortful, and the same holds in reverse for diminishment: a periphery whose salience has
been flattened stops issuing capture events at all, which is the mechanism this dissertation's dark
modes exploit.

The salience account has recently been extended to augmented reality specifically. Duan et al. (2022)
contribute the first dedicated AR saliency dataset (SARD) and demonstrate that the opacity of AR
content directly and predictably modulates where attention lands. This is empirical confirmation that
an alpha-blended overlay, the core primitive of this system's Tier-2 modes (Chapter 3), is an
attention instrument with a measurable dose–response character rather than a mere visual style.

\subsection{Cognitive load}

Where salience explains \textit{how} the periphery captures attention, Cognitive Load Theory explains \textit{why
that capture is costly}. Extraneous load, processing demanded by information irrelevant to the task,
competes for the same limited working-memory resources as the task itself. Buchner, Buntins and
Kerres (2022) review the evidence that AR can reduce, not only add, cognitive load when it is used to
structure or filter what the user must process. This supplies the theoretical justification for the entire
project: peripheral diminishment is an \textit{extraneous-load intervention}. It does not make the user
smarter or the task easier. It removes competing processing demands at the perceptual source. The
dissertation's primary hypothesis (H1, Chapter 6) is this theory operationalised: if peripheral
distractors impose a measurable performance cost on a focal task, a manipulation that suppresses
their salience should reduce that cost.

\subsection{The inherent tension}

This dissertation pairs two hypotheses, not one, because of a tension the literature above already
names. H1 says suppression reduces distraction cost. It says nothing about what suppression might
cost in return, and the same literature that motivates diminishment also names its central risk.
Attention capture by
peripheral events is not a design flaw of the human visual system. It is a safety mechanism.
Suppressing it trades focus against situational awareness, a tension documented empirically in the
DR literature (Murph et al., 2021; McLaughlin et al., 2025) and formalised in this
dissertation as an equivalence-tested safety hypothesis rather than an afterthought (H2b, Chapter 6).
This tension, which every section of this chapter returns to, is the reason the dissertation's
evaluation design refuses to measure benefit without simultaneously measuring cost.

\section{Saliency Modulation and Gaze Guidance in Mediated Views}

The techniques for steering attention in a mediated view form a spectrum from imperceptible to overt.
This section traverses it: subtle gaze direction (2.2.1), saliency modulation of video see-through
imagery, the direct ancestor of ColorPop (2.2.2), and overt cues with the comparative studies that
justify combining techniques (2.2.3).

\subsection{Subtle gaze direction}

The subtle gaze direction (SGD) family originates with Bailey, McNamara, Sudarsanam and Grimm (2009):
a brief luminance modulation ($\pm$9.5\%, $\sim$0.76$^{\circ}$ region, 10 Hz) presented in the viewer's periphery
attracts a saccade, and the cue is terminated the moment gaze moves toward it, so the viewer never
foveates the modulation and remains unaware of being steered. Luminance modulation outperformed
warm–cool chromatic modulation, an early instance of the luminance-dominance law of Section 2.1.
McNamara, Bailey and Grimm (2008) demonstrated task-level value: in visual search, subtle modulation
raised accuracy from 40.97\% to 56.25\%, statistically indistinguishable from an \textit{obvious} cue
(56.94\%), with zero of eighteen participants noticing the manipulation. Sridharan, Bailey, McNamara
and Grimm (2012) extended the paradigm to mammography training (accuracy 54.2\% $\rightarrow$ 64.5–68.8\%), and a
later SAP paper added automatic target prediction so the cue could terminate before foveal scrutiny
without an explicit eye-tracker trigger (Sridharan \& Bailey, 2015a). This is temporal termination
logic that informed this system's transition behaviour. A companion paper the same year extended the
paradigm to depth, guiding gaze in stereo virtual environments rather than flat images (Sridharan,
Pieszala \& Bailey, 2015b).

Grogorick and colleagues systematised the paradigm for immersive displays. In head-mounted 360$^{\circ}$
environments, Grogorick, Stengel, Eisemann and Magnor (2017) measured peripheral detection thresholds
that scale linearly with eccentricity ($s_{static}(e) = 0.0098e + 0.1009^{\circ}$; $s_{dynamic}(e) = 0.0076e
+ 0.289^{\circ}$) and showed subtle stimuli cutting hidden-target search from 36.07 s to 8.20 s. Their
five-way comparison in an immersive dome (N = 102; Grogorick, Albuquerque and Magnor, 2018; companion
article Grogorick et al., ACM TAP, 2018) is the most complete map of the subtlety–effectiveness
terrain: local magnification (\textit{ZoomRect}) achieved the best guidance-per-noticeability ratio
($\sim$30\% of target saccades within one second, and only 3 of 102 participants recognised it), while
\textit{SpatialBlur}, whole-periphery blur, was the least subtle and actively irritated 59 of 102
participants. Their sober conclusion, "no method remains completely unnoticed", together with the
blur-irritation finding, directly shaped this project. Whole-periphery blur is not among the modes
carried forward (Chapter 4), and the design accepts noticeability rather than chasing
imperceptibility. Related stimuli complete the palette. Stereo inverse brightness modulation
exploits binocular rivalry (Grogorick et al., 2020), noteworthy here because it can only ever be
implemented on content the application renders itself, never on the OS-owned passthrough layer
(Chapter 3). Flicker exploits the fovea–periphery difference in critical flicker fusion, visible
peripherally yet invisible foveally at 60–72 Hz (Waldin, Waldner and Viola, 2017), though marginal on
passthrough hardware whose camera pipeline runs well below those rates. Recent CHI work revisits
flicker as a practical AR guidance augmentation (Sutton et al., 2024), and its explicit
subtle-versus-overt framing is adopted in this chapter's organisation. Dichoptic colour cues
(different hue per eye) produce reliable pop-out in optical see-through AR, with Erickson, Bruder
and Welch's (2023) parametric
analysis of hue versus saturation versus value informing this project's choice to manipulate saturation
and luminance rather than hue. Outside the headset entirely, projector-based systems have guided gaze by
changing the apparent visual appearance of real surfaces, using a pixel-shift blur cheaper than
Gaussian filtering (Miyamoto, Koike and Amano, 2018). This is the
real-world analogue of passthrough shader diminishment, and the source of a performance idea (offset
sampling over kernel filtering) echoed in this system's mobile-GPU budget decisions.

Slow ramps constitute the second escape hatch from the noticeability trade-off: Hata, Koike and Sato
(2016) showed that blur introduced gradually enough stays below awareness threshold while still
biasing gaze, because the guidance threshold sits below the awareness threshold. This finding is the
perceptual justification for this system's formation animation: effects ramp in over seconds rather
than switching on (Chapter 4).

\subsection{Saliency modulation of see-through video}

Where SGD adds a transient stimulus, saliency \textit{modulation} re-grades the image itself, raising
conspicuity where attention should go and lowering it elsewhere. This is the direct technical
lineage of ColorPop.

The foundational demonstration is Veas, Mendez, Feiner and Schmalstieg (2011): using an Itti-style
conspicuity model, they raised salience in a focus region and suppressed it in context regions of
video imagery (channel order: luminance first, then red–green, then blue–yellow; algorithm in Mendez,
Feiner and Schmalstieg, 2010). Modulated video was statistically indistinguishable from unmodulated
video for naive viewers, yet produced faster first fixations on target regions (t(35) = 2.916,
p < .01) and improved recall of modulated content (0.19 $\rightarrow$ 0.31, p = .008). This remains the strongest
evidence that video-mediated reality can be re-graded imperceptibly with measurable attentional
consequences. And video mediation is exactly what a passthrough headset is. Kalkofen, Mendez and
Schmalstieg (2007) contribute the focus-and-context vocabulary this dissertation uses throughout:
the design discipline of keeping a focus region maximally legible while context is compressed but
not destroyed.

The modern anchor is Sutton, Langlotz, Plopski, Zollmann, Itoh and Regenbrecht (2022): saliency
modulation of the \textit{real world} viewed through AR glasses, implemented as saturation adjustment plus
sigmoidal contrast remapping ($\alpha$ = 10, $\beta$ = 0.5), deliberately avoiding hue shifts, which participants
found confusing. Modulation drew significantly more fixations to target regions ($\chi^2$ = 25.57,
p < .001), but, an honest and important negative, was \textit{not} imperceptible. Against an overt circle
cue, modulation's advantage was that it preserved natural scene exploration where explicit markers
"glue" gaze to themselves. This dissertation's ColorPop mode implements Sutton et al.'s exact
sigmoidal contrast recipe on its salient-colour channel (Chapter 4), inheriting the philosophy:
accept visibility, preserve exploration, re-grade rather than annotate. The critical difference is
platform. Sutton et al. built a bench-mounted optical see-through rig whose additive optics can only
add light. The present system runs on consumer video see-through hardware with full subtractive
control of its rendered layers (Section 2.5).

An eye-tracked comparison in cinematic VR speaks to which \textit{diminishment} channel works: comparing
area darkening, context-based darkening, and desaturation for steering attention in 360$^{\circ}$ video, area
darkening guided attention most effectively while desaturation alone had almost no guidance effect
(Wang, Gan \& Li, 2020). It is a single result from one small-venue study, not a consensus finding,
but it corroborates this project's mode hierarchy: the dark modes are the guidance workhorses, and
desaturation appears only as one ingredient inside ColorPop's compound re-grading, never as a mode of
its own.

\subsection{Overt cues, comparisons, and combinations}

At the overt end of the spectrum, the Attention Funnel (Biocca, Tang, Owen et al., 2006), a 3D
tunnel of AR rings leading the head to an out-of-view target, improved search speed by 22\% and
consistency by 65\% while reducing mental workload by 18\%, establishing that aggressive guidance
pays in speed what it costs in scene engagement. The corresponding failure mode is \textit{attentional
tunnelling}: excessive allocation of attention to guided content at the cost of neglecting the
surrounding environment. This is something handheld-AR work has since isolated experimentally as a function of
task demand (Syiem et al., 2021). This is the risk most relevant to this dissertation's Hard Dark
mode, and one of the reasons its evaluation measures peripheral awareness rather than assuming it.
Later cinematic-VR work found that overt guidance can also simply fail:
Nielsen et al. (2016) compared a diegetic firefly cue against forced body rotation in 360$^{\circ}$ film,
establishing head/gaze direction as the outcome measure for guidance in exactly the medium this
project's video test scene uses. Comparative studies fill in the middle ground: Renner and Pfeiffer
(2017) showed gaze-adaptive peripheral cues reducing search time in narrow-FOV AR, and their
follow-up work in complex AR environments found multi-modal (visual + audio) combinations reducing
search time and errors further (Renner and Pfeiffer, 2020). Marquardt et al. (2020) demonstrated
that audio-tactile guidance can match visual guidance in accuracy while \textit{improving} situational
awareness. This is a finding this dissertation flags for future work as the natural safety companion
to peripheral dimming: critical events in a dimmed periphery could be announced on the untouched
auditory channel. Hein, Bernhagen and Bullinger (2019) report that combining a coarse guidance
channel with a fine one significantly outperforms either alone, which supports this system's
compound design: coarse peripheral attenuation plus fine salience lift inside the focus window.
Salience computation has also been used in the service of \textit{placement} rather than guidance,
choosing where AR labels can sit without occluding what matters (Rakholia, Hegde and Hebbalaguppe,
2018). This dissertation notes it as a technique repurposable for saliency-weighted dimming,
attenuating hardest where the periphery is most distracting, as suggested too by Lu, Duh and Feiner's
(2012) finding that scene clutter should modulate cue intensity. Together with the Attention Funnel
above, these overt-cue studies establish the baseline this project defines itself against.
This system currently delegates focus-region selection to the user (trigger-painted windows) or to
pre-baked detections, by design (Chapter 4).

\section{Diminished Reality and Visual Noise Cancellation}

\subsection{Definitions and taxonomy}

Mori, Ikeda and Saito (2017) provide the canonical survey and taxonomy of diminished reality: the
family of techniques that \textit{conceal, eliminate, or see through} real objects in a mediated view. Their
survey establishes the field's legitimacy, that mediation of reality can subtract as well as add,
and its historical centre of gravity: most classical DR work concerns object removal (inpainting a
region so the object appears gone), which is computationally demanding and semantically targeted.
The present work sits in a different corner of their taxonomy: it does not remove objects but
\textit{attenuates regions}, trading semantic precision for robustness, generality, and real-time
feasibility on mobile hardware.

Cheng, Yin, Yan, Gugenheimer and Lindlbauer (2022) supply the field's most direct study of how users
\textit{want} reality diminished. Across two studies (N = 16 and N = 12, conducted inside VR simulations of
DR because no deployable DR hardware existed), they found participants preferred partial opacity
reduction over complete removal, and wanted \textit{context-preserving} diminishment: outlines or residual
structure indicating that something is there, even when its detail is suppressed. Both findings are
load-bearing for this dissertation: the preference for partial attenuation motivates the Soft Dark
mode's deliberately incomplete dimming (maximum opacity $\approx$ 0.75, never a blackout), and the
context-preservation finding motivates treating Hard Dark's total blackout as the \textit{extreme} of the
design space rather than its default. Equally load-bearing is Cheng et al.'s method: they could only
simulate DR inside VR. This dissertation runs on the real passthrough hardware their participants
could only imagine. The delta is developed in Section 2.8.

\subsection{Visual noise cancellation}

The most recent conceptual frame for this family of systems is \textit{visual noise cancellation} (VNC): by
analogy with acoustic noise cancellation, a head-worn display actively moderates the visual
environment, attenuating "noise" while passing "signal" (Hong et al., 2024). The analogy is productive
because it imports an architecture: acoustic ANC requires a
microphone (sensing the environment), a processing stage, and a speaker (re-emitting a modified
signal), precisely the camera $\rightarrow$ shader $\rightarrow$ display pipeline of a video-see-through headset. The VNC
paper is the closest conceptual match to this entire project, and this dissertation can be read as
the first \textit{software-defined VNC system on consumer video-passthrough hardware}: each mode is a VNC
profile, and the profiles differ in what they classify as noise and in how aggressively they
attenuate it. The dark modes treat everything peripheral as noise. SignPop treats everything that is
not a detected signal object as noise, so its rule is semantic rather than purely spatial. A follow-up line of work has begun evaluating VNC-style workspaces with sustained
attention and working-memory batteries in mixed-reality settings (see Section 2.7).

\subsection{DR for distraction suppression: the empirical record}

The empirical case that diminishment helps focus is recent but consistent. McLaughlin et al. (2025)
provide the closest published analogue to this dissertation's primary hypothesis: in two
within-subject experiments (STEM graduates and NASA Johnson Space Center employees), participants
performed a demanding assembly task under visual and auditory distraction, with DR-style attenuation
of the distractors applied either universally or context-sensitively. Universal attenuation, the
condition structurally closest to this project's peripheral vignette, produced fewer errors
($\beta = -0.37$, p = .021) and lower subjective workload than context-aware removal. Their analysis
(linear mixed-effects models over trial-level data) is adopted directly by this dissertation's
analysis plan. Critically, McLaughlin et al. also measured awareness of off-task objects and events,
finding the focus/awareness trade-off real but manageable. This is the empirical ancestor of this
work's H2b.

Lee and Kim's DiminishAR (2025) approaches the same goal from the object side: AR-camouflaging a
single distracting object (a smartphone) restored working-memory performance to a level statistically
comparable with physically removing it (N = 60). Murph et al. document the same benefit–risk
structure in a medical assembly context: DR lowered cognitive workload but risked situational
awareness loss (Murph et al., 2021). In later work they describe methods of training users to
overcome distraction under diminished reality, in which distractions are gradually reintroduced as
the trainee's tolerance grows (Murph, Richardson \& McLaughlin, 2022). That gradual-reintroduction
logic maps directly onto this system's temporal fade transitions (formation ramps and
motion-triggered suppression, Chapter 4). Richardson et al. address the complementary problem this
dissertation defers to future work: how to alert users to genuinely critical events inside a
diminished region (Richardson et al., 2021).

Two 2026 threads confirm both the demand for and the immaturity of DR-for-attention. A co-design
study with fifteen students with ADHD (DIS 2026) developed the concept of \textit{attentional diminishment}
through diary studies and workshops, blur, desaturation, and removal as user-configurable filters,
but built no working prototype (Exploring Diminished Reality for Attention Support, DIS 2026). A CHI
2026 study on Apple Vision Pro obscured specific individuals to alleviate social discomfort,
demonstrating overlay-based person-targeted diminishment on consumer hardware, though as targeted
occlusion rather than peripheral attention guidance (Obscuring Undesirable Individuals…, CHI 2026).
Both are formative or adjacent. Neither implements peripheral diminishment for attention on real
passthrough. Earlier applied work points the same direction. Yantaç, Corlu, Fjeld and Kunz (2015)
explored interactive DR walls filtering irrelevant visual information for individuals on the autism
spectrum, validating the psychological premise of environment filtering for populations for whom
distraction is most costly. FocalSpace (Yao, DeVincenzi, Pereira and Ishii, 2013) demonstrated the 2D
ancestor of this project's workstation scenario: synthetic background blur in video conferencing,
driven by depth and activity tracking, improving memory for meeting content and user preference.

\subsection{Hardware routes to diminishment}

Finally, two projects demonstrate the \textit{hardware} appetite for what this dissertation does in
software. IlluminatedFocus spatially defocuses regions of the real world using electrically tunable
lenses, reshaping perception optically (Ueda, Iwai and Sato, 2019). Zhang's Programmable Peripheral
Vision explores hardware and conceptual approaches to augmenting the peripheral visual field (Zhang,
2022). Both accept substantial optical complexity to achieve region-selective manipulation of
perceived reality. That the same capability now falls out of a shader on a €550 consumer headset is
the platform shift this dissertation documents in Section 2.5 and Chapter 3.

\section{Peripheral Vision, Foveated Degradation, and Comfort}

\subsection{What the periphery can lose}

Peripheral vision is not blurry central vision. It is a differently tuned system with its own
sensitivities (to motion and flicker above all) and its own tolerances. The review literature on
peripheral–foveal interactions (Stewart et al., 2020) establishes the parameters that bound
this project: how much degradation the periphery absorbs before scene perception destabilises, and
which manipulations the periphery notices most (contrast loss and motion change) versus least
(chromatic detail loss). Foveated rendering, degrading image quality with eccentricity to save
computation, has mapped this terrain thoroughly. The
state-of-the-art survey (Wang et al., 2022) catalogues the acceptable-degradation
mathematics this project reuses for its focus-to-periphery falloffs. Tursun et al. (2019) showed
that tolerable degradation is \textit{content-dependent}: a luminance-contrast-aware model permits far more
degradation in low-contrast regions, motivating this dissertation's recommendation (future work)
that diminishment intensity adapt to scene content rather than follow a fixed falloff.

Two perceptual findings sharpen the picture. Patney et al. (2016) established that peripheral blur is
detected primarily through the \textit{contrast reduction} it causes. Restoring local contrast after
blurring doubles the tolerable blur radius. Walton et al. (2021) push further with ventral metamers:
peripheral images that are structurally different but perceptually indistinguishable, degrading far
more gracefully than Gaussian blur. Both explain, mechanistically, the irritation finding of
Grogorick et al. (2018): naive whole-periphery blur announces itself through contrast collapse.
This chain of evidence is why the present system's dark
modes attenuate luminance (which the periphery tolerates) rather than blurring structure (which it
resents).

\subsection{Comfort limits of peripheral restriction}

A second literature approaches the periphery from the comfort side: field-of-view restriction as a
cybersickness countermeasure. Fernandes and Feiner (2016) showed a soft, dynamically applied FOV
vignette reduces discomfort in VR locomotion without measurable presence loss. This is the canonical
demonstration that peripheral occlusion is tolerable when it is soft and situationally applied.
Sustained use appears benign: Teixeira and Palmisano (2021) report sickness reductions holding over repeated
10-minute exposures with presence unaffected. But the manner of restriction matters greatly. Norouzi,
Bruder and Welch (2018) found that vignetting \textit{coupled to the user's own head motion}, strengthening
as the head turns, actually \textbf{increased} sickness relative to no vignette: a directly actionable
negative result that this system's motion policy inverts (effects fade \textit{out} during head motion and
restore during stability, Chapter 4), and the citation that justifies that inversion. Wu and Suma
Rosenberg (2022) showed that masking only the high-optic-flow periphery beats symmetric restriction
on presence: suppress the distracting \textit{signal}, not the region wholesale. Cao, Grandi and
Kopper (2021) found granulated rest frames (sparse static noise grains, 1–7$^{\circ}$, 25–75\% density)
outperform opaque black restrictors on search performance while preserving peripheral awareness,
with 15 of 20 participants preferring grains to blackout. Barhorst-Cates, Rand and Creem-Regehr
(2016) contribute the sobering datum on the far end: spatial learning survives restriction down to
10$^{\circ}$ FOV, but anxiety is elevated in \textit{all} restricted conditions and rises monotonically as the field
narrows. Waldner et al. (2014) map the comfort envelope for the flicker channel ($\approx \frac{1}{4}$ luminance range
below 2 Hz achieves $\geq$ 0.9 detection with low annoyance), parameters this project adopts for its
detection-highlighting pulse rather than anything in the 15–25 Hz photosensitivity band.

Collectively this literature draws the corridor the present system must fly through: peripheral
attenuation is tolerable and even beneficial (Fernandes \& Feiner), \textit{if} it is soft-edged,
decoupled from head motion (Norouzi), ideally selective about what it suppresses (Wu; Cao), and
mindful that total blackout carries a measurable anxiety cost (Barhorst-Cates) and an attention-
tunnelling risk (Biocca). Those constraints are visible in the final system as the dark modes'
graduated opacity ceilings, every mode's motion-triggered suppression, and the evaluation's insistence
on measuring peripheral awareness (H2b) and comfort (VRSQ, ratings) alongside benefit. One honest
gap remains, inherited rather than resolved: no study to date measures the comfort of peripheral
\textit{desaturation} specifically over sessions longer than ten minutes, and no head-to-head of blur
versus black vignette exists within a single design. This dissertation's end-of-session questionnaire collects
first subjective data on adjacent questions but does not close either gap.

\section{Platform Constraints and Prior Systems on Consumer Passthrough Hardware}

\subsection{The constraint landscape}

Everything in Sections 2.2–2.3 presumes the ability to manipulate what the user sees of reality. On
consumer hardware in 2026, that ability is narrowly rationed, and the rationing defines this
dissertation's technical contribution (Chapter 3 develops it as a three-tier design space).

On Meta Quest 3, the operating system owns the passthrough layer: applications cannot read its
pixels, and until recently could not modify them at all. Meta's Passthrough Styling API is the
exception that proves the rule: \verb|OVRPassthroughLayer| exposes brightness, contrast and saturation
controls, 3D colour look-up tables, and edge tinting, permitting \textit{global colour remapping} of the
entire passthrough layer, but no blur, no per-pixel spatial masks, and no world-locked regions
(Meta Passthrough Styling documentation). A global desaturation is expressible at this tier. A
peripheral or windowed effect is not. Surface-projected passthrough, which briefly allowed
per-surface styling, was deprecated at SDK v83. The complementary capability arrived in early 2025:
Passthrough Camera Access (PCA) grants applications the raw forward camera feed, full pixel
control, but of a stream markedly inferior to the system passthrough (lower resolution, $\sim$85–90$^{\circ}$
FOV against the system layer's $\sim$110$^{\circ}$, added latency, and no built-in stereo reprojection). One
published feasibility study so far quantifies what a native mixed-reality compositing engine
combining PCA with real-time segmentation could sustain on this hardware class. That is a
simulation-based estimate, not a measured on-device run, and its numbers double as a warning: it
projects 720p at 30 fps, with thermal throttling within five to ten minutes (Laghari et al., 2025,
arXiv:2509.18929). On Apple Vision Pro the rationing is stricter still: camera pixels require an
enterprise-only entitlement incompatible with consumer distribution, and the sole diminishment
primitive available to consumer apps is a global, non-spatial surroundings dimming (visionOS
documentation). Varjo's XR-3/XR-4 research headsets grant full camera control, yet no published study
was found that shader-processes their passthrough region-selectively for attention purposes.
Researchers with the capability have not asked this question, while consumer platforms that provoke
the question lack the capability.

\subsection{Prior systems and near-precedents}

Every constituent mechanic of this dissertation's architecture has a documented precedent, and none
of the precedents combine them. Camera-feed-on-geometry rendering with GPU effects exists in
community samples. QuestCameraKit demonstrates PCA-fed materials with blur and distortion shaders
applied to virtual panels and objects (xrdevrob/QuestCameraKit), but always as content \textit{within}
the scene, never as a full-surround treated periphery. A head-centred sphere whose shader-controlled
transparency reveals system passthrough exists as Meta's own "Passthrough Transitioning" motif, but
the sphere carries virtual content for scene fades, not a processed camera feed (Meta MR Motifs
documentation). Alpha punch-through to the passthrough underlay is a documented standard technique
("Passthrough Windows"), shipped commercially by Immersed's Passthrough Portals in 2022 so desk
workers could see their keyboards, but the hole is always surrounded by \textit{virtual} workspace, not by
re-rendered reality. A peripheral dimming vignette shipped in Quest OS v67 as "Theatre View", but it
explicitly does not operate over passthrough, only in immersive contexts. At the claim level, two
patents describe adjacent ideas, attenuating external stimuli while preserving spatial awareness
(US 12548271) and blurred external video feeds composited into VR (US 12524072), with no evidence
either was implemented.

A structured prior-art search conducted for this dissertation (5 July 2026; 26 web queries and 4
GitHub API queries across ISMAR, CHI, UIST, IEEE VR, DIS and arXiv 2022–2026, Meta developer
documentation, all 136 public forks of the PCA samples repository, XR trade press, and platform API
documentation; the full query protocol is archived in the project record) found \textbf{no publication,
product, or open-source project that composites system passthrough with a live shader-processed
camera feed in a single view, for any purpose}, and none that implements peripheral passthrough
manipulation for attention guidance on a standalone consumer headset. The nearest systems are the
five named above, each contributing one ingredient. The hybrid composite itself, native-quality
passthrough revealed through a world-locked window in an effect-processed camera-feed sphere,
exploiting the quality and field-of-view asymmetry between the two layers, appears to be novel, and
Chapter 3 presents it as such, scoped precisely to that composite rather than to "DR on passthrough"
in general.

\subsection{The perceptual ceiling of passthrough}

One further platform datum constrains not the system but its \textit{evaluation}: psychophysical
measurement shows that no current video see-through headset reaches normal human acuity at any
light level, with Quest 3 significantly worse in low light (Wang et al., 2026). Any evaluation task that
requires reading fine real-world detail through passthrough therefore confounds the manipulation
under test with the medium's acuity ceiling. This finding drives two decisions in Chapter 6: task
stimuli are either virtual or pilot-verified legible at large print, and room lighting is bright and
constant across sessions.

\section{Attention Guidance in Applied Contexts}

\subsection{Driving-adjacent contexts}

The densest applied literature on guiding attention under peripheral load is automotive. AR head-up
display research shows that salience-engineered cues, dynamic contour highlighting of hazards,
reduce inattentional blindness and improve reaction times, particularly in adverse weather (Zhu, Li
\& Liu, 2025). Controlled simulator studies corroborate: highly specific visual AR warnings improve
braking reaction time over generic alerts (Schwarz \& Fastenmeier, 2017; N = 88, within-subject), and
AR cueing of roadside hazards improves response time
\textit{without} degrading detection of non-target objects (Rusch et al., 2013; N = 27). The latter is
an early empirical
demonstration that guidance benefit and situational-awareness cost are separable, the exact
separation this dissertation's H2a/H2b pair formalises. The DR-specific studies of Murph et al.
(2021) and McLaughlin et al. (2025), reviewed in Section 2.3, translate the same
benefit–risk structure into diminishment terms.

A scoping statement is required here, because this dissertation uses driving \textit{footage} as one of its
stimulus contexts (Chapter 6). Per the claim-scoping rules fixed in the project's framing (Chapter
1), \textbf{no claims about driving, drivers, or road safety are made anywhere in this dissertation}. The
automotive literature above motivates the \textit{scenario aesthetics}, a dynamic, high-event-rate
monitoring stimulus with colour-coded semantic anchors (traffic signals), and supplies the
salience-engineering precedent for ColorPop's red/orange/green hierarchy. A seated participant
passively watching passenger-perspective footage licenses conclusions about attention mechanics
under dynamic stimulation, nothing more.

\subsection{Workstation and desk contexts}

The workstation scenario, this dissertation's primary evaluation context, has its clearest
ancestor in FocalSpace (Yao et al., 2013; Section 2.3), which established that suppressing background
distraction in a mediated view improves retention of foreground content. The recent VNC-adjacent
line of mixed-reality workspace studies (Section 2.3.2) brings the same question into headsets with
sustained-attention and working-memory batteries. What was missing until very recently was the
distractor-cost baseline in immersive settings. That has now been supplied. Peripheral visual
distractors in a virtual classroom significantly increase both commission errors (1.33 $\rightarrow$ 3.15,
p < .001) and omission errors (0.14 $\rightarrow$ 1.18, p < .001) on a Go/No-go continuous performance task
(N = 66, within-subject), while leaving mean reaction time unchanged (Ai et al., 2025). That single result
does double duty for this dissertation: it establishes that \textit{errors, not latency}, are the sensitive
currency of distraction cost (fixing the primary dependent variable of Chapter 6), and it quantifies
the effect the vignette must reduce for the primary hypothesis to be supported.

\subsection{Educational contexts}

The education literature motivated a third scenario, which the study does not run as a separate
block (Chapter 6). Its motivating findings stand: student attention drifts in instructional settings, and
gaze-informed interventions can steer it back. This is demonstrated by gaze-data visualisation systems for
educational VR (Rahman, Asish, Khokhar, Kulshreshth \& Borst, 2019) and, methodologically, by
gaze-contingent blur applied to dynamic sports stimuli in VR, whose evaluation apparatus (fixation
counts, area-of-interest dwell under central versus peripheral blur) validated that blur
manipulations shift gaze distribution measurably (Limballe, Kulpa, Vu, Mavromatis \& Bennett, 2022).
The present work retains these as design motivation
and evaluation precedent respectively, while declining to make confirmatory claims in an educational
context it no longer tests.

\section{Evaluating Attention Guidance in XR}

\subsection{Design and sample-size norms}

The evaluation norms of this field are within-subject designs with modest samples and many repeated
measures. Caine's census of CHI 2014 (2016) found the modal sample size across all user studies to
be twelve, with in-person studies smaller than remote ones. The nearest-neighbour DR studies ran
N = 16 and N = 12 (Cheng et al., 2022) and two within-subject experiments analysed with
mixed-effects models (McLaughlin et al., 2025). Sutton et al. (2022) ran N = 20 within-subject with
semi-randomised condition order. The methodological lesson codified in Chapter 6 is that small-N XR
studies derive their power from paired comparisons and trial-level repetition, not headcount, and
that trial-level linear mixed-effects models (McLaughlin's choice) extract more sensitivity from the
same sessions than aggregate t-tests. For non-normal aggregates the HCI norm has been the aligned
rank transform (Wobbrock et al., 2011), though recent critique of ART's error properties (Tsandilas
\& Casiez, 2024) motivates this dissertation's preference for mixed models with
Wilcoxon fall-backs, and Friedman tests with Bonferroni-corrected pairwise Wilcoxon comparisons for
Likert batteries (the Cheng pattern).

\subsection{Objective measures without eye tracking}

Quest 3 carries no eye tracker, so this dissertation's measures must come from behaviour and head
telemetry. The literature validates both channels. Reaction-time probe detection is standardised as
the Detection Response Task (ISO 17488:2016), with academic roots in Jahn, Oehme, Krems and Gelau
(2005): button-press latency to scheduled probe stimuli indexes spare attentional capacity. This
lineage motivated the driving-scene block of Chapter 6, though the task built and run there is
simpler than a full DRT implementation. Participants click on real traffic-light
and pedestrian markers as they appear in the driving footage, rather than on scripted abstract
probes, and onset time, position, hit, miss, and false alarm are all still machine-timed ground
truth. Error rates under distraction are the sensitive measure
for sustained-attention tasks (the Ai et al. result of Section 2.6.2: errors moved, latency did
not). Head orientation is a validated proxy for gaze at the eccentricities that matter here: Higgins
et al. (2022) formally validated head pose as a gaze surrogate for object-level attribution in VR,
and Sitzmann et al. (2018) showed with concurrent head and eye tracking in 360$^{\circ}$ environments that
fixations concentrate at low head velocity. Head yaw at rest is a reliable pointer to attention
beyond roughly 15$^{\circ}$ eccentricity, with eye tracking adding value only inside the central window.
Memory probes for attended versus unattended content (recognition of items placed at target versus
periphery) complete the eye-tracker-free toolkit.

\subsection{Subjective instruments and safety}

The standard subjective battery is stable across this literature. Workload: raw NASA-TLX, the
unweighted variant, which Hart's twenty-year retrospective (2006) endorses and Byers, Bittner and
Hill (1989) showed sacrifices no sensitivity against the pairwise-weighted original (Hart \&
Staveland, 1988). Simulator sickness: the SSQ (Kennedy et al., 1993) or, better suited to short
mixed-reality sessions, the nine-item VRSQ (Kim et al., 2018). Either must be administered before
\textit{and} after exposure, because post-only scores are uninterpretable without a baseline (Brown et al.,
2022). Of this battery the present study adopts the VRSQ, pre and post, and adds a single
end-of-session opinion questionnaire of its own. It administers no per-condition workload
instrument, and Chapter 6 records that decision and what it costs. Equivalence claims, this
dissertation's "no meaningful loss of peripheral awareness" hypothesis, require dedicated
machinery: the two-one-sided-tests procedure with a pre-registered margin (Lakens, 2017), which
converts a would-be null result into a falsifiable claim. Rothe, Buschek and Hußmann (2019) provide
the guidance-technique taxonomy used to situate all of the above. No canonical systematic review of
XR attention guidance covering 2020–2025 yet exists, a gap this chapter partially services but does
not claim to close.

\section{Synthesis: Research Gaps and Position}

\subsection{Gaps}

Read together, the eight bodies of literature above leave five identifiable gaps, three of which this
dissertation addresses in whole or part:

\begin{longtable}{|p{0.05\textwidth}|p{0.35\textwidth}|p{0.5\textwidth}|}
\hline
\textbf{\#} & \textbf{Gap} & \textbf{Status in this dissertation} \\
\hline
\endhead
1 & \textbf{No unified multi-mode DR system.} Guidance techniques and diminishment effects have been studied singly. No prior work combines multiple DR guidance modes in one configurable system for comparison within one design. & Addressed: ten modes built in one configurable system, spanning the design-space tiers of Chapter 3, of which two, SignPop and Hard Dark, are carried to confirmatory comparison within one study (Chapter 6). \\
\hline
2 & \textbf{No hybrid of system passthrough and app-processed camera feed in one view.} Colour-only global restyling of system passthrough exists (Meta Styling API). Full pixel control of an inferior camera feed exists (PCA). No prior work composites the two, and no prior work implements peripheral passthrough manipulation for attention on a standalone consumer HMD (documented search, Section 2.5.2). & Addressed: the window-as-absence hybrid architecture (Chapters 3–4), the dissertation's strongest and most precisely scoped novelty claim. \\
\hline
3 & \textbf{No temporal transition mechanics.} Gradual formation, easing, and motion-contingent suppression of DR effects have not been formally parameterised or studied, despite adjacent evidence that ramps evade noticeability (Hata et al., 2016) and that head-coupled strengthening harms comfort (Norouzi et al., 2018). & Partially addressed: formation ramps and inverted motion-coupling are implemented and parameterised (Chapter 4) and grounded in the cited evidence. Their isolated evaluation remains future work. \\
\hline
4 & \textbf{No user-controlled DR intensity.} User-initiated, graded release and re-application of diminishment is unstudied. & Not addressed, groundwork only (the trigger-painted window gives users spatial but not intensity control). Future work. \\
\hline
5 & \textbf{No dual-channel DR technique.} Desaturation and blur have been studied separately, never combined and parameterised as a compound diminishment channel in AR attention guidance. & Not addressed as a confirmatory question: the compound channel exists inside ColorPop, but the study design deliberately tests the \textit{unconfounded} dark modes confirmatorily and leaves channel decomposition to future work. \\
\hline
\end{longtable}

\subsection{Position against the nearest neighbours}

Two published studies stand closest to this dissertation, and the delta from each defines its
contribution. \textbf{Cheng et al. (2022)} established what users want from diminishment, but inside VR
simulations, because deployable DR did not exist. Their study is a perception-and-preference study
of imagined systems. \textbf{McLaughlin et al. (2025)} established that distractor attenuation improves
task performance, but likewise inside a virtual task environment, with the diminishment simulated
by the VR scene itself. Both are simulation studies \textit{of necessity}. This dissertation contributes
the piece both were missing: a working diminishment system on real consumer passthrough hardware
(Chapters 3–5), evaluated with a design that deliberately pairs a real-hardware block (the Hard Dark
workstation task, with physical distractors in the participant's actual room) against a
simulation block (the SignPop driving scene, a detection-gated peripheral re-grading, under oracle
perception) within the same participants
(Chapter 6). If the two blocks agree in direction, the field learns that the simulation methodology
of Cheng and McLaughlin predicts real-hardware outcomes. If they diverge, the divergence itself is a
finding about what simulation misses. Either way, the evaluation inherits the field's strongest
methodological norms: within-subject design, trial-level mixed models, pre-registered directional
and equivalence hypotheses, and a decision table that admits failure, assembled here, per the
review above, for the first time around a deployable diminished-reality attention system.

\section*{References}

\refitem{Bailey, R., McNamara, A., Sudarsanam, N., \& Grimm, C. (2009). Subtle Gaze Direction. \textit{ACM Transactions on Graphics}, 28(4), Art. 100. doi:10.1145/1559755.1559757}
\refitem{Barhorst-Cates, E. M., Rand, K. M., \& Creem-Regehr, S. H. (2016). The Effects of Restricted Peripheral Field-of-View on Spatial Learning while Navigating. \textit{PLoS ONE}, 11(10), e0163785.}
\refitem{Barreiros, C., Veas, E., \& Pammer-Schindler, V. (2016). Pre-attentive Features in Natural Augmented Reality Visualizations. \textit{2016 IEEE International Symposium on Mixed and Augmented Reality (ISMAR-Adjunct)}. doi:10.1109/ISMAR-Adjunct.2016.0043}
\refitem{Biocca, F., Tang, A., Owen, C., \& Xiao, F. (2006). Attention Funnel: Omnidirectional 3D Cursor for Mobile Augmented Reality Platforms. \textit{CHI 2006}. \url{https://ieeexplore.ieee.org/document/1579336/}}
\refitem{Syiem, B. V., Kelly, R. M., Goncalves, J., Velloso, E., \& Dingler, T. (2021). Impact of Task on Attentional Tunneling in Handheld Augmented Reality. \textit{CHI 2021}. doi:10.1145/3411764.3445580}
\refitem{Buchner, J., Buntins, K., \& Kerres, M. (2022). The Impact of Augmented Reality on Cognitive Load and Performance: A Systematic Review. \textit{Journal of Computer Assisted Learning}, 38(1), 285–303. doi:10.1111/jcal.12617}
\refitem{Byers, J. C., Bittner, A. C., \& Hill, S. G. (1989). Traditional and Raw Task Load Index (TLX) Correlations: Are Paired Comparisons Necessary? \textit{Advances in Industrial Ergonomics and Safety I}.}
\refitem{Caine, K. (2016). Local Standards for Sample Size at CHI. \textit{CHI 2016}. \url{https://dl.acm.org/doi/10.1145/2858036.2858498}}
\refitem{Cao, Z., Grandi, J., \& Kopper, R. (2021). Granulated Rest Frames Outperform Field of View Restrictors on Visual Search Performance. \textit{Frontiers in Virtual Reality}, 2:604889. doi:10.3389/frvir.2021.604889}
\refitem{Cheng, Y., Yin, H., Yan, Y., Gugenheimer, J., \& Lindlbauer, D. (2022). Towards Understanding Diminished Reality. \textit{CHI 2022}. doi:10.1145/3491102.3517452}
\refitem{Sridharan, S., Pieszala, J., \& Bailey, R. (2015b). Depth-based Subtle Gaze Guidance in Virtual Reality Environments. \textit{ACM SIGGRAPH Symposium on Applied Perception (SAP '15)}. doi:10.1145/2804408.2814187}
\refitem{Duan, H., et al. (2022). Saliency in Augmented Reality (SARD dataset). \textit{ACM Multimedia 2022}.}
\refitem{Erickson, A., Bruder, G., \& Welch, G. (2023). Analysis of the Saliency of Color-Based Dichoptic Cues in Optical See-Through Augmented Reality. \textit{IEEE Transactions on Visualization and Computer Graphics.} doi:10.1109/TVCG.2022.3195111}
\refitem{Exploring Diminished Reality for Attention Support: A Co-Design Study with Students with ADHD. \textit{DIS 2026}. \url{https://dl.acm.org/doi/10.1145/3800645.3813095}}
\refitem{Fernandes, A. S., \& Feiner, S. K. (2016). Combating VR Sickness through Subtle Dynamic Field-of-View Modification. \textit{IEEE 3DUI 2016}, 201–210. doi:10.1109/3DUI.2016.7460053}
\refitem{Sutton, J., Langlotz, T., Plopski, A., \& Hornbæk, K. (2024). Flicker Augmentations: Rapid Brightness Modulation for Real-World Visual Guidance using Augmented Reality. \textit{CHI 2024}. doi:10.1145/3613904.3642085}
\refitem{Wang, L., Shi, X., \& Liu, Y. (2022). Foveated Rendering: A State-of-the-Art Survey. arXiv:2211.07969.}
\refitem{Miyamoto, J., Koike, H., \& Amano, T. (2018). Gaze Navigation in the Real World by Changing Visual Appearance of Objects using Projector-Camera System. \textit{VRST 2018}. doi:10.1145/3281505.3281537}
\refitem{Grogorick, S., Stengel, M., Eisemann, E., \& Magnor, M. (2017). Subtle Gaze Guidance for Immersive Environments. \textit{ACM SAP 2017}. doi:10.1145/3119881.3119890}
\refitem{Grogorick, S., Albuquerque, G., \& Magnor, M. (2018). Comparing Unobtrusive Gaze Guiding Stimuli in Head-Mounted Displays. \textit{IEEE ICIP 2018}, 2805–2809. doi:10.1109/ICIP.2018.8451784. Companion: \textit{ACM TAP}, 15(4), 2018. doi:10.1145/3238303}
\refitem{Grogorick, S., et al. (2020). Stereo Inverse Brightness Modulation for Guidance in Dynamic Panorama Videos in VR. \textit{Computer Graphics Forum}, 39(2).}
\refitem{Hart, S. G. (2006). NASA-Task Load Index (NASA-TLX); 20 Years Later. \textit{HFES Annual Meeting}. \url{https://human-factors.arc.nasa.gov/groups/TLX/downloads/HFES_2006_Paper.pdf}}
\refitem{Hart, S. G., \& Staveland, L. E. (1988). Development of NASA-TLX. In \textit{Human Mental Workload}.}
\refitem{Hata, H., Koike, H., \& Sato, Y. (2016). Visual Guidance with Unnoticed Blur Effect. \textit{AVI 2016}, 28–35. doi:10.1145/2909132.2909254}
\refitem{Hein, P., Bernhagen, M., \& Bullinger, A. C. (2019). Improving Visual Attention Guiding by Differentiation between Fine and Coarse Navigation. \textit{2019 11th International Conference on Virtual Worlds and Games for Serious Applications (VS-Games)}. doi:10.1109/VS-Games.2019.8864539}
\refitem{Higgins, P., Barron, R., \& Matuszek, C. (2022). Head Pose as a Proxy for Gaze in Virtual Reality. \textit{VAM-HRI 2022}. \url{https://iral.cs.umbc.edu/Pubs/Higgins2022VAM-HRI.pdf}}
\refitem{Ai, X., Wang, Y., Wang, P., \& Wang, S. (2025). Impact of Visual Distractors in Virtual Reality Environments on Sustained Attention Behavioral Performance and EEG Characteristics. \textit{Frontiers in Human Neuroscience}. \url{https://pmc.ncbi.nlm.nih.gov/articles/PMC12698649/}}
\refitem{ISO 17488:2016. Road Vehicles — Detection-Response Task (DRT) for Assessing Attentional Effects of Cognitive Load in Driving.}
\refitem{Itti, L., Koch, C., \& Niebur, E. (1998). A Model of Saliency-Based Visual Attention for Rapid Scene Analysis. \textit{IEEE TPAMI}, 20(11).}
\refitem{Jahn, G., Oehme, A., Krems, J. F., \& Gelau, C. (2005). Peripheral Detection as a Workload Measure in Driving. \textit{Transportation Research Part F}, 8(3).}
\refitem{Kalkofen, D., Mendez, E., \& Schmalstieg, D. (2007). Interactive Focus and Context Visualization for Augmented Reality. \textit{ISMAR 2007}, 191–200. doi:10.1109/ISMAR.2007.4538846}
\refitem{Kennedy, R. S., Lane, N. E., Berbaum, K. S., \& Lilienthal, M. G. (1993). Simulator Sickness Questionnaire. \textit{International Journal of Aviation Psychology}, 3(3).}
\refitem{Kim, H. K., Park, J., Choi, Y., \& Choe, M. (2018). Virtual Reality Sickness Questionnaire (VRSQ). \textit{Applied Ergonomics}, 69, 66–73. \url{https://www.sciencedirect.com/science/article/abs/pii/S000368701730282X}}
\refitem{Lakens, D. (2017). Equivalence Tests: A Practical Primer. \textit{Social Psychological and Personality Science}, 8(4). doi:10.1177/1948550617697177}
\refitem{Lee, J., \& Kim, S. (2025). DiminishAR. \textit{CHI 2025}. doi:10.1145/3706598.3713415 (preprint: arXiv:2403.03875)}
\refitem{Lu, W., Duh, B.-L. H., \& Feiner, S. (2012). Subtle Cueing for Visual Search in Augmented Reality. \textit{2012 IEEE International Symposium on Mixed and Augmented Reality (ISMAR)}, 161–166. doi:10.1109/ISMAR.2012.6402553}
\refitem{Tursun, O. T., et al. (2019). Luminance-Contrast-Aware Foveated Rendering. \textit{ACM Transactions on Graphics}, 38(4). doi:10.1145/3306346.3322985}
\refitem{Marquardt, A., Trepkowski, C., Eibich, T. D., Maiero, J., Kruijff, E., \& Schöning, J. (2020). Comparing Non-Visual and Visual Guidance Methods for Narrow Field of View Augmented Reality Displays. \textit{IEEE TVCG}, 26(12), 3389–3401. doi:10.1109/tvcg.2020.3023605}
\refitem{McLaughlin, A. C., et al. (2025). Cognitive Aid Design Using Diminished Reality to Support Selective Attention by Reducing Distraction. \textit{Human Factors}, 67(9), 937–961. doi:10.1177/00187208251325169}
\refitem{McNamara, A., Bailey, R., \& Grimm, C. (2008). Improving Search Task Performance Using Subtle Gaze Direction. \textit{APGV 2008}, 51–56.}
\refitem{Mendez, E., Feiner, S., \& Schmalstieg, D. (2010). Focus and Context in Mixed Reality by Modulating First Order Salient Features. \textit{Smart Graphics 2010}, LNCS 6133, 232–243.}
\refitem{Meta Platforms. Passthrough Styling / Passthrough Windows / MR Motifs: Passthrough Transitioning / Passthrough Camera Access. Meta Horizon developer documentation. \url{https://developers.meta.com/horizon/documentation/unity/}}
\refitem{Mori, S., Ikeda, S., \& Saito, H. (2017). A Survey of Diminished Reality. \textit{IPSJ Transactions on Computer Vision and Applications}, 9:17. doi:10.1186/s41074-017-0028-1}
\refitem{Murph, I., McDonald, M., Richardson, K., Wilkinson, M., Robertson, S., Karunakaran, A., Gandy Coleman, M., Byrne, V., \& McLaughlin, A. C. (2021). Diminishing Reality: Potential Benefits and Risks. \textit{Proceedings of the Human Factors and Ergonomics Society Annual Meeting}, 65(1), 164–168. doi:10.1177/1071181321651103}
\refitem{Murph, I., Richardson, K., \& McLaughlin, A. C. (2022). Methods of Training to Overcome Distraction Via Diminished Reality. \textit{Proceedings of the Human Factors and Ergonomics Society Annual Meeting}, 66(1), 1844–1848. doi:10.1177/1071181322661134}
\refitem{Laghari, M. K., Shaikh, A. A., Khan, F., \& Siddiqui, A. G. (2025). Native Mixed Reality Compositing on Meta Quest 3: A Quantitative Feasibility Study of ARM-Based SoCs and Thermal Headroom. arXiv:2509.18929.}
\refitem{Nielsen, L. T., et al. (2016). Missing the Point: An Exploration of How to Guide Users' Attention During Cinematic Virtual Reality. \textit{VRST 2016}. doi:10.1145/2993369.2993405}
\refitem{Norouzi, N., Bruder, G., \& Welch, G. (2018). Assessing Vignetting as a Means to Reduce VR Sickness During Amplified Head Rotations. \textit{ACM SAP 2018}. doi:10.1145/3225153.3225162}
\refitem{Obscuring Undesirable Individuals to Alleviate Social Discomfort Using Diminished Reality. \textit{CHI 2026}. \url{https://dl.acm.org/doi/10.1145/3772318.3790918}}
\refitem{Patents US 12548271 B2 (Apple Inc., granted 2026-02-10) — "Attention control in multi-user environments" (title confirmed verbatim via Google Patents); US 12524072 B2 (Samsung Electronics, granted 2026-01-13) — "Providing a pass-through view of a real-world environment for a virtual reality headset for a user interaction with real world objects", paraphrased here as "blurred external video feed in VR". Both numbers confirmed valid and correctly assigned; claim-level adjacency only, no implementation evidence in either case.}
\refitem{Patney, A., et al. (2016). Towards Foveated Rendering for Gaze-Tracked Virtual Reality. \textit{ACM Transactions on Graphics}, 35(6). doi:10.1145/2980179.2980246}
\refitem{Wang, J., Ping, S., Xu, K., Li, Y., \& Liang, H.-N. (2026). The Perceptual Gap between Video See-Through Displays and Natural Human Vision. arXiv:2601.02805.}
\refitem{Zhang, Q. (2022). Programmable Peripheral Vision: Augment/Reshape Human Visual Perception. \textit{CHI '22 Extended Abstracts.} doi:10.1145/3491101.3503821}
\refitem{Rahman, Y., Asish, S. M., Khokhar, A., Kulshreshth, A. K., \& Borst, C. W. (2019). Gaze Data Visualizations for Educational VR Applications. \textit{ACM Symposium on Spatial User Interaction (SUI '19)}. doi:10.1145/3357251.3358752}
\refitem{Renner, P., \& Pfeiffer, T. (2017). Attention Guiding Techniques using Peripheral Vision and Eye Tracking. \textit{IEEE 3DUI 2017}.}
\refitem{Renner, P., \& Pfeiffer, T. (2020). AR-glasses-based Attention Guiding for Complex Environments: Requirements, Classification and Evaluation. \textit{PETRA '20.} doi:10.1145/3389189.3389198}
\refitem{Richardson, K., McLaughlin, A. C., McDonald, M., \& Crowson, A. (2021). The Effects of Diminished Reality on the Detection of and Response to Notifications. \textit{Proceedings of the Human Factors and Ergonomics Society Annual Meeting}, 65(1), 159–163. doi:10.1177/1071181321651236}
\refitem{Limballe, A., Kulpa, R., Vu, A., Mavromatis, M., \& Bennett, S. J. (2022). Virtual Reality Boxing: Gaze-Contingent Manipulation of Stimulus Properties Using Blur. \textit{Frontiers in Psychology}, 13:902043. doi:10.3389/fpsyg.2022.902043}
\refitem{Rothe, S., Buschek, D., \& Hußmann, H. (2019). Guidance in Cinematic Virtual Reality — Taxonomy, Research Status and Challenges. \textit{Multimodal Technologies and Interaction}, 3(1), 19.}
\refitem{Rakholia, N., Hegde, S., \& Hebbalaguppe, R. (2018). Where to Place: A Real-Time Visual Saliency Based Label Placement for Augmented Reality Applications. \textit{IEEE ICIP 2018.} doi:10.1109/ICIP.2018.8451052}
\refitem{Sitzmann, V., Serrano, A., Pavel, A., Agrawala, M., Gutierrez, D., Masia, B., \& Wetzstein, G. (2018). Saliency in VR: How Do People Explore Virtual Environments? \textit{IEEE TVCG}, 24(4), 1633–1642. (arXiv:1612.04335)}
\refitem{Sridharan, S., Bailey, R., McNamara, A., \& Grimm, C. (2012). Subtle Gaze Manipulation for Improved Mammography Training. \textit{ETRA 2012}, 75–82. doi:10.1145/2168556.2168568}
\refitem{Brown, P., Spronck, P., \& Powell, W. (2022). The Simulator Sickness Questionnaire, and the Erroneous Zero Baseline Assumption. \textit{Frontiers in Virtual Reality}, 3:945800. doi:10.3389/frvir.2022.945800}
\refitem{Stewart, E. E. M., Valsecchi, M., \& Schütz, A. C. (2020). A Review of Interactions between Peripheral and Foveal Vision. \textit{Journal of Vision}, 20(12):2. doi:10.1167/jov.20.12.2}
\refitem{Sridharan, S., \& Bailey, R. (2015a). Automatic Target Prediction and Subtle Gaze Direction for Improved Spatial Information Recall. \textit{ACM SAP 2015}, 99–106. doi:10.1145/2804408.2804415}
\refitem{Sutton, J., Langlotz, T., Plopski, A., Zollmann, S., Itoh, Y., \& Regenbrecht, H. (2022). Look over there! Investigating Saliency Modulation for Visual Guidance with AR Glasses. \textit{UIST 2022}, Art. 81. doi:10.1145/3526113.3545633}
\refitem{Teixeira, J., \& Palmisano, S. (2021). Effects of Dynamic Field-of-View Restriction on Cybersickness and Presence in HMD-Based Virtual Reality. \textit{Virtual Reality}, 25(2), 433–445. doi:10.1007/s10055-020-00466-2}
\refitem{Tsandilas, T., \& Casiez, G. (2024). The Illusory Promise of the Aligned Rank Transform. \textit{Journal of Visualization and Interaction} (under review). \url{https://statransform.github.io/jovi/}}
\refitem{Ueda, T., Iwai, D., \& Sato, K. (2019). IlluminatedFocus: Vision Augmentation using Spatial Defocusing. \textit{SIGGRAPH Asia 2019 Emerging Technologies.} doi:10.1145/3355049.3360530}
\refitem{Veas, E., Mendez, E., Feiner, S., \& Schmalstieg, D. (2011). Directing Attention and Influencing Memory with Visual Saliency Modulation. \textit{CHI 2011}, 1471–1480. doi:10.1145/1978942.1979158}
\refitem{Hong, J., Langlotz, T., Sutton, J., \& Regenbrecht, H. (2024). Visual Noise Cancellation: Exploring Visual Discomfort and Opportunities for Vision Augmentations. \textit{ACM Transactions on Computer-Human Interaction.} doi:10.1145/3634699}
\refitem{Waldin, N., Waldner, M., \& Viola, I. (2017). Flicker Observer Effect: Guiding Attention through High Frequency Flicker in Images. \textit{Computer Graphics Forum}, 36(2), 467–476. doi:10.1111/cgf.13141}
\refitem{Waldner, M., Le Muzic, M., Bernhard, M., Purgathofer, W., \& Viola, I. (2014). Attractive Flicker — Guiding Attention in Dynamic Narrative Visualizations. \textit{IEEE TVCG}, 20(12), 2456–2465. doi:10.1109/TVCG.2014.2346352}
\refitem{Walton, D. R., Kuffner Dos Anjos, R., Friston, S., Swapp, D., Akşit, K., Steed, A., \& Ritschel, T. (2021). Beyond Blur: Real-Time Ventral Metamers for Foveated Rendering. \textit{ACM Transactions on Graphics (SIGGRAPH 2021)}, 40(4), 1–14. doi:10.1145/3450626.3459943}
\refitem{Wang, G., Gan, Q., \& Li, Y. (2020). Research on Attention-guiding Methods in Cinematic Virtual Reality Based on Eye Tracking Analysis. \textit{2020 International Conference on Innovation Design and Digital Technology (ICIDDT)}. doi:10.1109/ICIDDT52279.2020.00020}
\refitem{Schwarz, F., \& Fastenmeier, W. (2017). Augmented Reality Warnings in Vehicles: Effects of Modality and Specificity on Effectiveness. \textit{Accident Analysis \& Prevention}, 101, 55–66. doi:10.1016/j.aap.2017.01.019}
\refitem{Rusch, M. L., Schall, M. C. Jr., Gavin, P., Lee, J. D., Dawson, J. D., Vecera, S., \& Rizzo, M. (2013). Directing driver attention with augmented reality cues. \textit{Transportation Research Part F: Traffic Psychology and Behaviour, 16}, 127–137. doi:10.1016/j.trf.2012.08.007}
\refitem{Wobbrock, J. O., Findlater, L., Gergle, D., \& Higgins, J. J. (2011). The Aligned Rank Transform for Nonparametric Factorial Analyses Using Only ANOVA Procedures. \textit{Proceedings of CHI 2011}, 143–146. doi:10.1145/1978942.1978963}
\refitem{Wu, F., \& Suma Rosenberg, E. (2022). Adaptive Field-of-view Restriction: Limiting Optical Flow to Mitigate Cybersickness in Virtual Reality. \textit{VRST 2022}. doi:10.1145/3562939.3565611 (\textbf{corrected 2026-08-05}: the title previously given here — "Asymmetric FOV Restriction Using Optic Flow" — does not match this DOI. That title belongs to a different Wu/Suma Rosenberg paper, \textit{Don't Block the Ground: Reducing Discomfort in Virtual Reality with an Asymmetric Field-of-View Restrictor} (SUI 2021, doi:10.1145/3485279.3485284). The DOI was correct; the title has been fixed to match it.)}
\refitem{Coviello, R. [xrdevrob] (2025). \textit{QuestCameraKit} [software]. GitHub. \url{https://github.com/xrdevrob/QuestCameraKit}}
\refitem{Yantaç, A. E., Corlu, D., Fjeld, M., \& Kunz, A. (2015). Exploring Diminished Reality (DR) Spaces to Augment the Attention of Individuals with Autism. \textit{2015 IEEE International Symposium on Mixed and Augmented Reality Workshops}, 68–73. doi:10.1109/ISMARW.2015.21}
\refitem{Yao, L., DeVincenzi, A., Pereira, A., \& Ishii, H. (2013). FocalSpace: Multimodal Activity Tracking, Synthetic Blur and Adaptive Presentation for Video Conferencing. \textit{ACM SUI 2013}. doi:10.1145/2491367.2491377}
\refitem{Zhu, Q., Li, J., \& Liu, Y. (2025). Visual Saliency Design for AR-HUD Navigation in Extreme Weather: Reducing Inattentional Blindness. \textit{IEEE Access}, 13, 137613–137622. doi:10.1109/ACCESS.2025.3588576}
